\DocAppendixTitle{Core Derivations}
\label{app:derivations}

\paragraph{Gibbs variational identity.}
Let \(P_\beta\) be the target in Equation~\eqref{eq:policy-target}. For any
\(Q\ll\mu\),
\begin{align}
  \KL(Q\|P_\beta)
  &= \E_Q\left[
       \log\frac{dQ}{d\mu}
       - \beta\ExpectedReturn_H(\policy)
       + \log Z_\beta
     \right] \\
  &= \log Z_\beta
     - \left(
       \beta\E_Q[\ExpectedReturn_H(\policy)]
       - \KL(Q\|\mu)
     \right).
\end{align}
Non-negativity yields Equation~\eqref{eq:gibbs-variational}, with equality at
\(Q=P_\beta\).

\paragraph{Why a rollout is valid in the direct objective.}
For \(\policy\sim q_\theta\) and
\(\tau\sim p(\cdot\mid\policy)\), iterated expectation gives
\[
  \E_{\policy,\tau}[\Return_H(\tau)]
  = \E_{\policy\sim q_\theta}
    [\ExpectedReturn_H(\policy)].
\]
No exponential or logarithm is applied between the inner and outer
expectations.  Applying the score-function identity to
Equation~\eqref{eq:correct-elbo}, and dropping terms whose expected score is
zero, gives Equation~\eqref{eq:score-gradient}.

\paragraph{Local detailed balance.}
At inverse temperature \(b\), write
\[
  \widetilde P_b(\policy)\propto
  \mu(\policy)\exp(b\widetilde J(\policy)).
\]
The Metropolis--Hastings ratio is
\[
  \frac{\widetilde P_b(\policy')q(\policy\mid\policy')}
       {\widetilde P_b(\policy)q(\policy'\mid\policy)}
  =
  \exp\!\left[b\{\widetilde J(\policy')-
                   \widetilde J(\policy)\}\right]
  \frac{\mu(\policy')}{\mu(\policy)}
  \frac{q_s(a\mid a')}{q_s(a'\mid a)}.
\]
The uniform decision-state selection cancels, producing
Equation~\eqref{eq:local-mh}.  Full support makes every one-coordinate change
possible, and a finite sequence of such changes connects any two tabular
policies.

\paragraph{Replica-swap ratio.}
Before a swap, the factors for replicas \(i\) and \(j\) are proportional to
\(
  \mu(\policy_i)\exp(b_i\widetilde J(\policy_i))
  \mu(\policy_j)\exp(b_j\widetilde J(\policy_j))
\).
After exchanging the two policies, the reference-measure factors cancel and
the new-to-old density ratio is
\[
  \exp\!\left[(b_j-b_i)
    \{\widetilde J(\policy_i)-\widetilde J(\policy_j)\}\right],
\]
which gives Equation~\eqref{eq:swap-mh}.  Alternating disjoint adjacent pairs
changes scheduling efficiency but not this invariant product distribution.

\paragraph{Fixed-tape convergence.}
For every fixed policy, the strong law gives
\(\widehat J_K(\policy)\to J(\policy)\) almost surely.  Because the tabular
policy space is finite, convergence holds simultaneously for all policies.
Continuity of exponentiation and finite normalization then imply
\(P_{\beta,K}(\policy\mid\Xi_K)\to P_\beta(\policy)\) for every policy.
This argument justifies a convergence study in \(K\); it does not make a
finite tape bank exact for \(P_\beta\).

\paragraph{A nonnegative unbiased policy weight from bounded rollouts.}
Suppose a simulator return satisfies \(G_\policy\geq a\) almost surely, where
the finite-horizon reward specification supplies the known constant \(a\).
For \(\lambda>0\), draw \(N\sim\operatorname{Poisson}(\lambda)\) and, conditional
on \(N\), independent returns \(G_{\policy,1},\ldots,G_{\policy,N}\).  Define
\begin{equation}
  \widehat W_{\lambda}(\policy)
  = \exp(\beta a)
    \prod_{i=1}^{N}
      \left(1+\frac{\beta(G_{\policy,i}-a)}{\lambda}\right).
  \label{eq:poisson-policy-weight}
\end{equation}
Every factor is nonnegative.  Conditioning first on \(N\), then using the
probability generating function of a Poisson variable, gives
\begin{align}
  \E[\widehat W_{\lambda}(\policy)]
  &= \exp(\beta a)\,
     \E_N\!\left[
       \left(1+\frac{\beta(J(\policy)-a)}{\lambda}\right)^N
     \right] \\
  &= \exp(\beta a)\exp(\beta(J(\policy)-a))
   = \exp(\beta J(\policy)).
\end{align}
Thus a pseudo-marginal Metropolis--Hastings chain on the policy and estimator
randomness has the exact policy target as its marginal
\citep{andrieuRoberts2009,jacobThiery2015}.  The identity is not automatically
an efficient algorithm.  Its expected rollout count is \(\lambda\), and direct
calculation gives
\begin{equation}
  \frac{\E[\widehat W_{\lambda}(\policy)^2]}
       {\exp(2\beta J(\policy))}-1
  = \exp\!\left(
      \frac{\beta^2\E[(G_\policy-a)^2]}{\lambda}
    \right)-1.
  \label{eq:poisson-policy-weight-rv}
\end{equation}
Bounded relative variance requires rollout cost scaling with
\(\beta^2\E[(G_\policy-a)^2]\), potentially prohibitive.

\section{Exact one-decision diagnostic}
\label{app:noise-objective}

Consider one state, one decision, and two policies. The safe action returns
$c$ surely; the variable action returns $+1$ or $-1$ with equal probability.
There are no state revisits, exploration difficulties, resampling ancestry,
or policy-representation differences. Set $\beta=1$ and use a uniform prior.
The expected-return target and exponential-trajectory target respectively give
\begin{align}
 P_{\rm mean}(\text{variable})
   &=\frac{1}{e^c+1}, \\
 P_{\rm trajectory}(\text{variable})
   &=\frac{\cosh(1)}{e^c+\cosh(1)}.
 \label{eq:noise-targets}
\end{align}
When $c=0$, both policies have mean zero, so the first probability is $1/2$;
the second is approximately $0.606776$. When $c=0.2$, the safe policy has the
higher mean, yet the trajectory target still assigns the variable policy
approximately $0.558180$ mass. The expected-return target instead assigns
the safe policy approximately $0.549834$ mass. Preference for variability is
therefore a substantive objective choice, not an arithmetic error.

For categorical proposal $q_\theta$ with logits $\theta=(\theta_0,\theta_1)$,
the direct objective is
\[
 \mathcal L(\theta)=q_\theta(0)c+H(q_\theta)-\log2.
\]
At uniform initialization its gradient is $(c/4,-c/4)$. A single independent
return per policy supplies an unbiased score-function estimate of this
gradient, although individual updates remain noisy.

\begin{samepage}
To contrast an evidence objective, draw $N$ independent actions and returns,
and form
\[
 W_i=\frac{e^{G_i}}{2q_\theta(A_i)},\qquad
 \ell_N=\log\!\left(\frac1N\sum_iW_i\right),\qquad
 F_N(\theta)=\E[\ell_N].
\]
\end{samepage}
Writing $\bar W_i=W_i/\sum_jW_j$ and
$s_i=\nabla_\theta\log q_\theta(A_i)$, differentiation gives
\begin{equation}
 \nabla_\theta F_N
 =\E\!\left[-\sum_i\bar W_i s_i+\ell_N\sum_i s_i\right].
 \label{eq:noise-evidence-gradient}
\end{equation}
The first term differentiates the importance weights; the second accounts
for sampling actions from the parameterized proposal. This formula is not
the direct expected-return gradient. It is also not asserted to equal the
gradient of a variational KL to the trajectory target at every finite $N$.

Each independent draw has three possible outcomes, with probabilities
$q_\theta(0)$, $q_\theta(1)/2$, and $q_\theta(1)/2$. Enumerating all $3^N$
combinations therefore computes Equation~\eqref{eq:noise-evidence-gradient}
without Monte Carlo error. At uniform logits, the variable-action component is:
\begin{center}
\begin{tabular}{rrrr}
\toprule
$c$ & $N$ & Evidence gradient & Direct ELBO gradient \\
\midrule
0   & 1 & 0.000000  & 0.000000 \\
0   & 2 & 0.054223  & 0.000000 \\
0   & 4 & 0.045058  & 0.000000 \\
0.2 & 1 & -0.050000 & -0.050000 \\
0.2 & 2 & 0.023860  & -0.050000 \\
0.2 & 4 & 0.028336  & -0.050000 \\
\bottomrule
\end{tabular}
\end{center}
\begin{samepage}
A positive component initially increases preference for the variable action.
For $N=1$, taking the expectation of $G-\log2-\log q_\theta(A)$ recovers
the direct ELBO on this instance. That special identity does not equate the
objectives for larger $N$, or justify multi-state entropy approximations.
\end{samepage}

The executable check compares the one-step differentiable implementation with
this enumeration for $N=1,2,4$, both values of $c$, and uniform or nonuniform
logits: twelve cases and 96 saved numeric arrays. Independent inspection
verifies the outcomes, masses, gradients, source identity and compute record.
Native float32 discrepancies remain below the predeclared $3\times10^{-6}$
absolute tolerance. Cross-platform re-summation of derived float64 quantities
is checked separately at a 64-machine-epsilon bound; the native tolerance is
not enlarged. The largest averaged native-gradient discrepancy is
$3.70\times10^{-8}$.

This check isolates the one-step objective distinction. It does not prove
that either update learns a maze better, explain a full-system performance
gap, or establish the scalability of either inference engine. Those questions
require separately controlled, replicated experiments.
